\documentclass[11pt,a4paper]{article}

\usepackage[utf8]{inputenc}
\usepackage[T1]{fontenc}
\usepackage[margin=2.5cm]{geometry}
\usepackage{amsmath}
\usepackage{amssymb}
\usepackage{xcolor}
\usepackage{enumitem}
\usepackage{float}
\usepackage{hyperref}
\usepackage{parskip}
\usepackage{booktabs}
\usepackage{framed}
\usepackage{tcolorbox}
\usepackage{etoolbox}  % for \ifstrequal etc.

% ---------------------------------------------------------------
% Colors
% ---------------------------------------------------------------
\definecolor{revblue}{RGB}{30,80,160}
\definecolor{respgreen}{RGB}{20,120,60}
\definecolor{respyellow}{RGB}{180,140,20}
\definecolor{secpurple}{RGB}{120,40,140}
\definecolor{lightblue}{RGB}{224,236,255}
\definecolor{lightgreen}{RGB}{224,248,230}
\definecolor{lightyellow}{RGB}{255,250,205}
\definecolor{lightgray}{RGB}{245,245,245}

% ---------------------------------------------------------------
% tcolorbox styles
% ---------------------------------------------------------------
\tcbuselibrary{skins,breakable}

\tcbset{
  revstyle/.style={
    enhanced, breakable,
    colback=lightblue, colframe=revblue,
    boxrule=1.2pt, arc=2pt,
    left=8pt, right=8pt, top=6pt, bottom=6pt,
    before skip=8pt, after skip=4pt,
    colbacktitle=revblue, coltitle=white,
    fonttitle=\bfseries,
  },
  respstyle/.style={
    enhanced, breakable,
    colback=lightgreen, colframe=respgreen,
    boxrule=1.2pt, arc=2pt,
    left=8pt, right=8pt, top=6pt, bottom=6pt,
    before skip=4pt, after skip=10pt,
    colbacktitle=respgreen, coltitle=white,
    fonttitle=\bfseries,
  },
  unfinishedrespstyle/.style={
    enhanced, breakable,
    colback=lightyellow, colframe=respyellow,
    boxrule=1.2pt, arc=2pt,
    left=8pt, right=8pt, top=6pt, bottom=6pt,
    before skip=4pt, after skip=10pt,
    colbacktitle=respyellow, coltitle=white,
    fonttitle=\bfseries,
  },
}

% ---------------------------------------------------------------
% Numbering system
%   \newreviewer{label}  -- resets the per-reviewer counter,
%                           stores the reviewer label (e.g. "R1")
%   \revcomment{title}{body} -- increments both counters,
%                               prepends "Rn.k |" to the title
%   \response{body}          -- echoes same number in the title
% ---------------------------------------------------------------
\newcounter{revtotal}   % absolute comment counter across all reviewers
\newcounter{revlocal}   % per-reviewer comment counter
\newcommand{\currentreviewer}{}  % stores short label, e.g. R1

\newcommand{\newreviewer}[1]{%
  \setcounter{revlocal}{0}%
  \renewcommand{\currentreviewer}{#1}%
}

\newcommand{\revcomment}[2]{%
  \stepcounter{revtotal}%
  \stepcounter{revlocal}%
  \edef\commentlabel{\currentreviewer.\arabic{revlocal}}%
  \begin{tcolorbox}[revstyle,
      title={\commentlabel\ $|$\ #1}]
    #2
  \end{tcolorbox}%
  % store label for \response to reuse
  \edef\lastcommentlabel{\commentlabel}%
}

\newcommand{\response}[1]{%
  \begin{tcolorbox}[respstyle,
      title={Response to \lastcommentlabel}]
    #1
  \end{tcolorbox}%
}

\newcommand{\unfinishedresponse}[1]{%
  \begin{tcolorbox}[unfinishedrespstyle,
      title={Unfinished Response to \lastcommentlabel}]
    #1
  \end{tcolorbox}%
}

\newcommand{\seealso}[1]{\textit{\textcolor{gray}{(see also our responses to: #1)}}}

% ---------------------------------------------------------------
% Section formatting (manual, no titlesec needed)
% ---------------------------------------------------------------
\renewcommand{\section}[1]{%
  \vspace{1.2em}%
  {\large\bfseries\color{secpurple} #1}\\[-0.4em]%
  {\color{secpurple}\rule{\linewidth}{0.8pt}}%
  \vspace{0.4em}%
}

\newcommand{\subsect}[1]{%
  \vspace{0.8em}%
  {\normalsize\bfseries\color{revblue} #1}%
  \vspace{0.3em}\par%
}

% ---------------------------------------------------------------
\begin{document}

% ---------------------------------------------------------------
%  TITLE
% ---------------------------------------------------------------
\begin{center}
  {\LARGE\bfseries Response to Reviewers}\\[0.8em]
  {\large CardamomOT: a mechanistic optimal transport-based framework\\
  for gene regulatory network inference, trajectory reconstruction\\
  and generative modeling}\\[0.5em]
  {\normalsize Manuscript Number: PCOMPBIOL-D-26-00826}\\[0.3em]
  {\normalsize Yann Maugé and Elias Ventre}\\[0.3em]
  {\normalsize \today}
\end{center}

\bigskip

We thank the Editor, Section Editor, and all four Reviewers for their careful and
constructive reading of our manuscript. Their comments have helped us identify important
improvements in clarity, validation, and presentation. Below we address each point
individually. Changes to the manuscript are highlighted in the revised version
(\textit{Revised Manuscript with Track Changes}).

% ===============================================================
\section{Response to the Editor --- Journal Requirements}
% ===============================================================

\newreviewer{Ed}

\revcomment{Journal Requirement --- Manuscript source file}{%
Please upload your manuscript file as a .doc, .docx, .rtf or .tex (LaTeX Source File).}
\response{%
We have uploaded the manuscript source as a \texttt{.tex} file under the item type
``LaTeX Source File'', with the corresponding PDF as ``Manuscript''.}

\revcomment{Journal Requirement --- Figure files}{%
Please upload all main figures as separate Figure files in .tif or .eps format.}
\response{%
All main figures have been exported as separate \texttt{.eps} files and added to the submission.}

\revcomment{Journal Requirement --- Supporting Information}{%
Supplementary Figures and Tables are included in the manuscript file. Please remove them
and upload them with the file type ``Supporting Information'', with a legend listed in the
manuscript.}
\response{%
Supplementary figures and tables have been removed from the main manuscript file and
uploaded as separate Supporting Information files. Legends for each Supporting Information
item have been added to the manuscript after the reference list.}

\revcomment{Journal Requirement --- Financial Disclosure}{%
Please amend your detailed Financial Disclosure statement (full sentences, initials
alongside each funding source, role of funders).}
\response{%
The financial disclosure statement has been amended as follows: ``This work was supported
by an ENS de Lyon CDSN PhD fellowship awarded to YM. The funders had no role in study
design, data collection and analysis, decision to publish, or preparation of the
manuscript.''}

\section{Additional changes}

\noindent We detail here a modification that has been done during the revision process, independently on what the reviewers had noticed:\par

\textbf{Extension of \textsc{CardamomOT} to multiple explicit stimulus nodes.} 
The Schiebinger reprogramming dataset involves two 
biochemically distinct perturbation phases: doxycycline-induced OSKM expression 
(days~0--8) and serum/LIF culture conditions (days~0--18). In the previous 
version of the algorithm, \textsc{CardamomOT} was limited to a single constant 
stimulus node, which conflated these successive perturbations into an averaged 
effect whose inferred targets lacked clear biological interpretability. We have 
now extended the algorithm to support multiple simultaneous stimulus nodes, each 
with a user-defined activity window, and re-analyzed this dataset accordingly. 
Two stimulus nodes were introduced---\emph{Dox} and \emph{Serum}---yielding 
biologically distinct and more interpretable regulatory programs for each 
perturbation (Figure S3, E,G). The \emph{Dox} node captures the 
absence of activation of late pluripotency markers such as \emph{Dppa5a} and 
the epithelial gene \emph{Cldn7}, consistent with the early reprogramming phase 
preceding the mesenchymal-to-epithelial transition. The \emph{Serum} node 
drives repression of \emph{Lox}, a pro-EMT enzyme whose suppression is known 
to facilitate mesenchymal-to-epithelial transition during reprogramming, as well 
as genes associated with inflammatory signaling (\emph{Ereg}, \emph{Tnfrsf12a}). 
This decomposition demonstrates that \textsc{CardamomOT} can now disentangle the 
contributions of successive or co-occurring perturbations, improving both the 
biological interpretability of inferred GRNs and the internal consistency of the 
model. 
This re-analysis affects all Schiebinger-derived results (Figures 2, 6, 7, S3–S5, S10 and Table S3), and we have verified that it leaves every conclusion of the original submission unchanged — in particular the predicted enhancement of iPSC-like conversion under Obox6 and Zfp42 overexpression, with an unchanged magnitude of effect. Obox6 additionally now ranks among the top candidate hub regulators (Table S3). The manuscript and supplementary figures have been updated accordingly.\par


% ===============================================================
\section{Response to Reviewer \#1}
% ===============================================================

\newreviewer{R1}

We thank Reviewer~1 for a thorough and rigorous evaluation of the manuscript.
We address each major and minor comment in turn.

\subsect{Major Comments}

\revcomment{Benchmark fairness should be clarified}{%
The synthetic data appear to be generated from a model closely related to the assumptions
of CardamomOT, which may favour the proposed method. The authors should either add
benchmarks with model misspecification or explicitly frame the synthetic benchmarking as
validation under model-matched conditions.}
\response{%
This is a fair and important point. The simulated benchmarks are generated from the same
mechanistic model underlying CardamomOT (via the HARISSA package), constituting a
\textit{model-recovery test} rather than a fully general benchmark. We have clarified this
framing explicitly in Section~2.1: the simulated benchmarks establish
best-case performance when model assumptions are satisfied, and all competing methods are
evaluated under identical conditions.

Benchmarks under model misspecification are an important direction for future work. We partly addressed this by testing the robustness of CardamomOT to technical noise and irregular temporal sampling in our new analysis \seealso{R2.6, R4.5}}

\revcomment{Uncertainty in inferred GRNs should be reported}{%
Experimental GRNs are presented as point estimates. The authors should report edge
stability or confidence, for example using bootstrapping, subsampling, or repeated random
initializations.}
\response{%
We agree that edge uncertainty is important for biological interpretation. We have added:
\begin{itemize}
  \item A supplementary figure (Figure~S9) showing GRN stability across repeated runs with different random seeds for the experimental datasets.
  \item Expanded Discussion (Scalability and identifiability) on run-to-run variability as a diagnostic for identifiability and guidance for use.
\end{itemize}}

\revcomment{Perturbation predictions should be more cautiously interpreted}{%
The distinction between predictions supported by independent experimental evidence and
purely model-generated hypotheses should be made clearer. Effect sizes and variability
across repeated simulations would be more informative than only reporting chi-squared
tests.}
\response{%
We have revised Section~2.6 to improve distinction between predictions consistent with prior experimental findings or molecular functions (Dnmt3a OV in the Semrau dataset; Obox6 and Zfp42 OV in the Schiebinger dataset), and model-generated hypotheses awaiting experimental validation (CHGA/STMN2 perturbations in the Kameneva dataset).

Figure~7 now includes error bars indicating variability across 5 independent repeated simulations per perturbation alongside the $\chi^2$ test results, which provide a more interpretable measure of biological relevance than $p$-values alone.}

\revcomment{Claims moderation}{%
Several claims should distinguish validation on simulated data from consistency checks on
experimental data. Statements such as ``accurately recovers velocity fields'', ``accurately
infers GRNs'', and ``reliable GRN structures'' should be moderated when referring to
experimental datasets.}
\response{%
We have revised the Abstract, Author summary, and Results sections throughout to consistently
distinguish:
\begin{itemize}
  \item \textit{Simulated data}: quantitative claims are retained where ground truth is
        available (e.g.\ ``cosine similarity $>0.8$ on benchmark networks'').
  \item \textit{Experimental data}: replaced with consistency-based language (e.g.\ ``the calibrated model reproduces the observed expression dynamics'', ``the inferred
        GRN is consistent with known regulatory roles'').
\end{itemize}
The Abstract now reads: ``We validate our framework on both in silico and experimental datasets, demonstrating computational scalability and consistently improved performance over state-of-the-art methods in both GRN and trajectory reconstruction on simulated datasets, and, on experimental datasets, reconstruction of cellular trajectories, velocity fields and latent protein levels that are mutually consistent, together with GRN structures consistent with known biology.''}

\subsect{Minor Comments}

\revcomment{Inconsistent terminology: ``gene regulation network'' vs ``gene regulatory network''}{%
The authors use both terms. Please choose one.}
\response{%
We have standardised the terminology throughout to ``gene regulatory network'' (GRN). The
keyword list has been updated accordingly.}

\revcomment{Hybrid stochastic-deterministic notation in Eq.~(2)}{%
The two-state model is described as stochastic for both $M_i(t)$ and $P_i(t)$, but Eq.~(2)
combines stochastic dynamics for $M_i$ production with mean-field ODEs for degradation and
$P_i$. Please revise or explain this hybrid notation.}
\response{%
The model in Eq.~(2) is indeed a \textit{hybrid stochastic-deterministic} (piecewise
deterministic) model: mRNA production occurs via a Poisson process driven by burst events
(stochastic), while degradation of $M_i$ and the dynamics of $P_i$ are deterministic ODEs.
This is the standard formulation of the two-state telegraph model in the bursty regime,
consistent with the framework of Piecewise Deterministic Markov Processes (PDMPs).

We have added a clarifying sentence in Section~1.1 making the hybrid nature of Eq.~(2)
explicit, and added a notation remark distinguishing stochastic jump terms (written with
$\rightarrow$) from deterministic flow terms (ODEs).}

\revcomment{Define notation $\varepsilon(k_{\mathrm{off}}/s_0)$}{%
E.g.\ does this denote an exponential distribution?}
\response{%
Yes --- $\mathcal{E}(k_{\mathrm{off}}/s_0)$ denotes an \textit{exponential distribution}
with rate parameter $k_{\mathrm{off}}/s_0$ (mean burst size $s_0/k_{\mathrm{off}}$). We have added an explicit definition at first use in Section~1.1.}

\revcomment{Define the negative binomial parameterisation in Eq.~(3)}{%
Different conventions use different meanings for the second parameter.}
\response{%
We agree, and we have added an explicit definition at first use in Section~1.1.
We write $\mathrm{NB}(a,\,b)$ for the Gamma--Poisson mixture defined by
$\lambda \sim \mathrm{Gamma}(\text{shape } a,\ \text{rate } b)$ and $M \mid \lambda \sim \mathcal{P}(\lambda)$,
so that in Eq.~(3)
\[
  a = \frac{k^{\theta}_{\mathrm{on},i}(P(t))}{d_{0,i}}
  \quad\text{is the \emph{shape} parameter (burst frequency relative to the mRNA degradation rate)},
\]
\[
  b = \frac{k_{\mathrm{off},i}}{s_{0,i}}
  \quad\text{is the \emph{rate} parameter (inverse mean burst size)}.
\]
This convention gives $\mathbb{E}[M_i] = a/b = \big(k^{\theta}_{\mathrm{on},i}(P(t))/d_{0,i}\big)\,\big(s_{0,i}/k_{\mathrm{off},i}\big)$
and a Fano factor $1 + s_{0,i}/k_{\mathrm{off},i}$, i.e.\ the classical result that mean expression is the
product of burst frequency and burst size, while overdispersion is controlled by burst size alone.
For readers using the size--probability convention, this is equivalent to a negative binomial with
size $a$ and probability $b/(1+b)$. We have added this equivalence explicitly in the text so that the
parameterisation is unambiguous under either convention.}

\revcomment{Additional peer-reviewed source for the NB approximation}{%
Could the authors provide an additional peer-reviewed source, in addition to the PhD thesis
currently cited [31]?}
\response{%
We have added peer-reviewed references supporting the NB approximation, including
Zong et al.\ (2010, \textit{Mol.\ Syst.\ Biol.}) and Gr\"{u}n et al.\ (2014, \textit{Nature Methods}), cited
alongside~[31] at first use.}

\revcomment{Algorithm~1: should there be $N-1$ OT problems?}{%
Apologies if my indexing is wrong --- should there be $N-1$ OT problems rather than $N$?}
\response{%
The reviewer is correct. For $N$ timepoints, one OT problem is solved per consecutive pair
$(t_j, t_{j+1})$, yielding $N-1$ problems. Algorithm~1 has been corrected.}

\revcomment{Abbreviations: define at first use, use consistently}{%
E.g.\ ``Reference Fitting (RF). Reference Fitting'' appears redundant.}
\response{%
We have reviewed all abbreviations and ensured they are defined at first use and used
consistently thereafter. The redundant occurrence has been corrected.}

\revcomment{State number of simulations and meaning of error bars in captions}{%
Please state number of independent simulations and what error bars/bounds represent in
figure captions.}
\response{%
We modified figure captions to include the number of independent simulations and the
nature of error bars (standard error of the mean or standard deviation, as appropriate) consistently.}

\revcomment{Figure~7: detail on reconstructed cell counts for $\chi^2$ tests}{%
E.g.\ are all $p < 10^{-3}$ because you ran 1000 simulations?}
\response{%
We thank the reviewer for this remark. We agree that p-values computed from the raw simulated cell counts would be strongly influenced by sample size. We have therefore changed the way the $\chi^2$ statistic is computed in the current version to decouple it from the number of simulated cells. The test compares the cell-type distributions of Sim perturb and Sim single. Rather than using the raw simulated cell counts, which range from a few thousand cells for Semrau and Kameneva to $\sim$250{,}000 for Schiebinger, we rescale the compared proportions to a fixed pseudo-count of 100 cells per cell type before constructing the contingency table. This provides a common scale for comparing the magnitude of compositional shifts across datasets, independently of their original sample sizes. We have made explicit in both Section 2.6 and the caption of Figure 7 that this rescaled statistic is used as a comparative measure of compositional separation rather than as a formal inferential test on the underlying biological system, since the pseudo-count is fixed by convention rather than derived from the data. Effect variability is reported separately as error bars over $N=5$ independent stochastic replicates. Under this revised analysis, CHGA KO alone no longer produces a shift clearly distinguishable from the run-to-run variability of the simulations, and we have adjusted the corresponding statement in Section 2.6 accordingly.
}


\revcomment{``Digital twin'' terminology may not be appropriate}{%
The model does not appear to involve iterative updating between the model and the
biological system.}
\response{%
We agree. ``Digital twin'' implies bidirectional feedback between a model and a physical
system, which is not the case here. We have replaced it with `` the 
calibrated model functions as an in-silico representation of the biological system'', which more accurately describes the CardamomOT output.}

\revcomment{Dataset accession numbers missing from Data and Code Availability}{%
Please ensure that all experimental datasets have clear accession numbers, not only a
GitHub link.}
\response{%
Explicit accession numbers have been added to the Data and Code Availability section:
\begin{itemize}
  \item Semrau et al.\ (2017): GEO accession GSE79578.
  \item Kameneva et al.\ (2021): GEO accession GSE147821.
  \item Schiebinger et al.\ (2019): GEO accession GSE106340 and \url{https://singlecell.broadinstitute.org/single_cell/study/SCP295/}
\end{itemize}
These have also been added to the GitHub README.}

% ===============================================================
\section{Response to Reviewer \#2}
% ===============================================================

\newreviewer{R2}

We thank Reviewer~2 for their careful attention to presentation and methodological
assumptions.

\revcomment{Figure~1: mixed vector/raster rendering}{%
Figure~1 could be cleaned up: it appears to be a jarring mix of vector text and raster
equations of varying resolution; panel~B seems to have vector axes but low-resolution
raster points.}
\response{%
We have regenerated Figure~1 so that all elements (text, equations, data points) are
rendered as vectors. The figure has been re-exported and uploaded.}

\revcomment{Figure~1 panels B and E: unclear relationship}{%
I am unsure what panels B and E represent in relation to each other. Panel~E (predictive
trajectories) seems to be a subset of the data points of panel~B (experimental plot), when
I would expect them to be switched.}
\response{%
Panels B and E contrast the \textit{input} (B: raw experimental snapshots) with the model
\textit{output} (E: simulated trajectories from the calibrated model). We have revised the
figure layout to make this contrast visually explicit, with 
\begin{itemize}
  \item Clear ``CardamomOT input'' / ``CardamomOT output'' labels.
  \item Redrawing the ''cell'' points to explicit resimulation at step E and avoid confusion due to the previous similarity of the cells.
\end{itemize}
}

\revcomment{Figure~3: ``Trees'' box border clipped}{%
In Figure~3, the ``Trees'' box border has been clipped.}
\response{%
This has been corrected in the revised Figure~3.}

\revcomment{Validity of piecewise-linear interpolation outside the adiabatic regime}{%
A considerable emphasis is placed on the model's extension to a non-linear regime. Is
there a flaw in the methodology from discretising the trajectories into linear segments using
the piecewise constant parameter $\alpha$ (Appendix A)? The justification seems to rest on the
asymptotic regime where the protein timescale $\ll$ mRNA degradation, but what about
outside this regime?}
\response{%
The reviewer's question rightly pointed to a step that deserved to be stated explicitly rather
than left implicit, and we have restated it more carefully than in the original submission.

Eq.~(15) in Appendix~A is not a numerical discretization of a nonlinear ODE. Within an interval,
we replace the protein-dependent burst frequency $k^{\theta}_{\mathrm{on},i}(P(t))$ by a
piecewise-constant driving term taking the two endpoint basin modes, with a single switch;
under that substitution the protein ODE~(5) is linear and Eq.~(15) is its exact closed-form
solution, so no additional discretization error arises from the formula itself. We should
however be explicit that the substitution \emph{is} the approximation, and it has two parts: the
regulatory feedback is frozen within the interval, and a single switch per gene is assumed to
describe the interval.

What we would emphasize is the scope of this piecewise-constant treatment, which we had not made
clear enough. It applies only to the construction of candidate protein paths between two
observed timepoints. The nonlinear GRN-driven structure is retained wherever the network is
estimated or used: in the OT cost of Eq.~(8), in the GRN regression of Eq.~(10) and the basin
refinement of Eq.~(11), in the kinetic recalibration of Section~1.2.3, where the full nonlinear
field is integrated with an adaptive-step solver and no piecewise-constant substitution, and in
the generative simulations, which sample the complete PDMP~(2). The linearity of Eq.~(15) is a
local device for generating transport candidates, not a linearization of the inferred dynamics;
the contrast with a linear reference process such as RF's Ornstein--Uhlenbeck model therefore
stands.

Both parts of the approximation are justified by the metastable, bursty and adiabatic regime
already required for the mechanistic model of Section~1.1 (see also our responses to: R4.8,
R4.3): within a basin the driving mode is close to constant, and transitions between
well-separated basins are rare and fast. Outside this regime, the discrete-basin representation
underlying the entire method --- the Negative Binomial mixture and its basin labels, not just
the protein interpolation --- ceases to describe the data before the interpolation itself becomes
the binding issue: the two approximations share a regime of validity and fail together, not
independently. We have rewritten the \emph{Regime of validity} paragraph of Appendix~A
accordingly.}



\revcomment{$n_{\mathrm{pas}}$: choice in practice and effect of coarse sampling}{%
The linear-spaced sample set (line 904) raises questions: (a) how is $n_{\mathrm{pas}}$ determined in
practice? (b) What effect does a coarse sample space have on the misidentification of
candidate cells and therefore the trajectory?}
\response{%
(a) $n_{\mathrm{pas}}$ is the number of grid points used to search for the switching time
$\alpha_{c,i}$ within an inter-timepoint interval (Appendix~A). The default is
$n_{\mathrm{pas}}=25$, and the implementation automatically increases it so that at least one
grid point falls within each time unit of $\Delta t = t_{j+1}-t_j$; this default was used,
unmodified, for every benchmark and experimental dataset in this article, so in practice this
is not a parameter users need to tune.

(b) The effect of a coarser grid is bounded and, in the regimes relevant to this work,
negligible. Given a fixed switching gene, Eq.~(15) varies smoothly in
$\alpha$, so a coarser grid only introduces a smooth interpolation error in the reconstructed
protein value rather than a discontinuous misassignment. More importantly, the search only
needs enough resolution to determine on which side of the mode-crossing time the interval
falls (Appendix~A, item~3). Finally, $\alpha_{c,i}$ is not estimated once: it is refined at every iteration of
the EM loop jointly with the GRN, so any residual error from a coarse initial grid tends to be
corrected over subsequent iterations rather than propagating to the final trajectories. We
have added a remark to this effect in Appendix~A. We would expect this parameter to become
relevant only in a pathologically coarse regime ($n_{\mathrm{pas}}$ of a few grid points,
effectively removing timing resolution within the interval), which is well below both the
default and the automatically enforced minimum. 

To illustrate this, we performed a sensitivity analysis under $n_{\mathrm{pas}}$ refinement as proposed by the reviewer (Figure~S11). Performance plateaus from $n_{\mathrm{pas}} \approx 4$ onwards; only the extreme value
$n_{\mathrm{pas}} = 2$, which removes essentially all timing resolution within an interval,
degrades accuracy, consistent with the pathologically coarse regime anticipated above. The
default of $25$ thus sits far inside the plateau.}


\revcomment{Potential bias from timestep choice in benchmarks}{%
The choice of timestep influences the switching parameter sample set. (a) More detail on
the choices of times and $n_{\mathrm{pas}}$ used in benchmark generation. (b) Sensitivity
analysis of Figure~3 results under timestep and $n_{\mathrm{pas}}$ refinement.}
\response{%
\begin{enumerate}[label=(\alph*)]
  \item The benchmark timepoints (0--96~h) and $n_{\mathrm{pas}} = 25$ were fixed prior to
        any comparisons and are identical for all methods. Competing methods do not use
        $n_{\mathrm{pas}}$ at all, so there is no differential tuning. These choices are
        fully documented in the code repository.
  \item We have added a sensitivity analysis (Figure 4H) showing GRN inference AUPR
  under increasing variability in the inter-timepoint intervals: for each replicate, the gaps
  between consecutive timepoints are drawn from a Dirichlet distribution with the endpoints
  of the time course held fixed, so that $\eta = 0$ recovers equispaced snapshots and
  $\eta = 1$ corresponds to uniformly sampled ones. Randomising the gaps
  at fixed endpoints necessarily produces replicates in which some consecutive timepoints
  are widely separated — precisely the configurations in which the switching-parameter
  sample set has to resolve a transition spread over a long interval.
  It also mimics the situation faced experimentally, where the number of snapshots is fixed
  by budget and the times at which the informative transitions occur are not known in
  advance, so that some intervals inevitably straddle them. AUPR is stable across the range of $\eta$, with slightly increased variability at high $\eta$ indicating that the benchmark results do not depend much on a favourable
  placement of the sampling times.

  We should be explicit about what this analysis does not vary: the total number of
  snapshots is held constant, so it probes the placement of the timepoints rather than a
  uniform refinement or coarsening of the time course as a whole. We would expect the
  latter to be governed by the temporal-resolution limitation already discussed in
  Section 3 — namely that intervals containing several mode switches per gene fall outside
  the single-switch assumption of Appendix A — rather than by any interaction with
  $n_{\mathrm{pas}}$, since $n_{\mathrm{pas}}$ is automatically increased with $\Delta t$ so
  that grid resolution is preserved as intervals lengthen.
        
\end{enumerate}}

\revcomment{Nuanced discussion of the method hierarchy in Figure~3}{%
Should the reader conclude that CardamomOT is currently the best method and CARDAMOM the
second best? It is difficult to discern what is responsible for this hierarchy --- optimal
transport minimising squared-Euclidean cost appears worse than neglecting OT mechanics
altogether.}
\response{%
We agree this ordering is easy to misread, and we have added a dedicated paragraph to
Section~2.3 to make the explanation explicit. The hierarchy in Figure~3A should not be read as a
ranking of methods in general, and in particular not as evidence that OT-based trajectory
inference is counterproductive; it reflects how well each method's dynamical and statistical
assumptions match the process that generated the benchmark data.

CARDAMOM never reconstructs trajectories, but its per-timepoint statistical model is the
Negative Binomial approximation derived from the very mechanistic model used to simulate the
benchmarks (Section~2.1), so it is, in this specific sense, correctly specified for these data
and can recover accurate parameter estimates without tracking cells across timepoints. RF does
reconstruct trajectories, but under a linear Ornstein--Uhlenbeck reference process that neither
exploits this NB structure nor captures the nonlinear, GRN-driven dynamics of the benchmark
networks; this mismatch degrades both the inferred trajectories and the GRN subsequently built
from them. CardamomOT is designed precisely to remove this trade-off, by pairing explicit
trajectory reconstruction with the same nonlinear mechanistic reference process instead of a
misspecified linear one --- which is why it outperforms both. The general statement is therefore
that the value of an OT step depends on how well the reference process matches the true
dynamics, not on the presence or absence of trajectory inference per se.

We emphasize that the \texttt{noloop}
ablation of Figure~4A--E, which performs a single round of OT on the burst modes without the iterative
refinement loop, supports this reading directly: its accuracy is lower than that of the full method, substantially so on BN8, yet it remains above RF on the same networks (compare Figure~4 and Figure~3A).}


% ===============================================================
\section{Response to Reviewer \#3}
% ===============================================================

\newreviewer{R3}

We thank Reviewer~3 for their positive assessment and constructive suggestions.

\revcomment{General assessment and recommendations}{%
The manuscript is well organised and the methods are described in sufficient detail to allow
reproduction. The results on simulated benchmarks are convincing. The application to
experimental datasets is promising, although the support there is somewhat more qualitative.
It would strengthen the manuscript to include more quantitative evaluation on experimental data
where feasible, and a brief discussion of sensitivity to key hyperparameters and model
assumptions. Overall, a strong contribution; recommend acceptance after minor revision.}
\response{%
We thank Reviewer~3 for this positive evaluation. In response to the suggestions:
\begin{itemize}
  \item More quantitative evaluation on experimental data: we have added stability analysis
        for experimental GRNs \seealso{R1.2}, perturbation effect sizes with variability
        across runs \seealso{R1.3}, and PHATE embeddings as an additional
        dimensionality reduction method \seealso{R4.1}.
  \item Hyperparameter sensitivity and model assumptions: We discuss sensitivity to protein degradation rates in more details \seealso{R4.3}, as well as robustness towards data modulations \seealso{R2.6, R4.5} and provide guidance for users in the discussion.
\end{itemize}}

% ===============================================================
\section{Response to Reviewer \#4}
% ===============================================================

\newreviewer{R4}

We thank Reviewer~4 for a detailed and constructive review. Several major suggestions have
led to meaningful improvements in the manuscript.

\subsect{Major Comments}

\revcomment{Robustness of low-dimensional embedding validation}{%
Several validation analyses rely on visual comparison in UMAP space. The authors should
repeat key visualizations using an additional trajectory-preserving method such as PHATE.}
\response{%
We have added PHATE embeddings as supplementary figures (Figure~S10) for all three
experimental datasets, showing both reference and simulated data. The PHATE embeddings
confirm the conclusions drawn from UMAP: the CardamomOT simulated data preserves the
global trajectory structure and temporal ordering of the reference data. PHATE embeddings
were computed from the concatenated reference and simulated datasets, with each plotted
separately, consistent with our UMAP methodology.}

\revcomment{Gene subset selection and possible selection bias}{%
The manuscript would benefit from a clearer discussion of how gene subsetting may affect
the inferred networks, and whether CardamomOT can be run on substantially larger gene
sets.}
\response{%
We have added explanations regarding these points in Section~3:
\begin{itemize}
  \item \textbf{Gene selection introduces biases, and is itself prior information.}
  We have added a dedicated paragraph to the Discussion (``Gene selection and its influence on
  inferred networks''). It first discusses how the choice of gene set shapes
  the inferred network: the influence of an omitted regulator is redistributed among the retained
  genes, so that inferred edges may summarise indirect regulation through unobserved
  intermediates; and selecting genes from known markers or previously characterised modules, as
  for the Semrau and Kameneva datasets, biases the network toward well-characterised regulatory
  axes. We now state explicitly that the agreement between our inferred hub regulators and those
  reported in the original studies is therefore a consistency check on a pre-selected gene set
  rather than an unbiased recovery of the underlying network, and we note that purely statistical
  selection criteria are not exempt either.

  We then argue that the gene set should be recognised as a form of prior information in its own
  right, on the same footing as prior constraints on individual edges, and give a practical
  recommendation: use prior biological knowledge of the system, or the hypotheses motivating the
  analysis, to designate a core set of nodes of interest, then assemble the regulatory context of
  those nodes semi-automatically by combining curated interaction databases with data-driven
  discovery --- for instance a genome-wide scalable method, used to obtain a
  coarse-grained approximation of the global network from which the neighbourhood of the core set
  is retained. Automating this construction is the object of ongoing work in our team.

  \item \textbf{Scalability, and the trade-off with prior structure.} Table~S2 reports runtimes
  up to 100 genes; as discussed in the Discussion (``Scalability and identifiability''), the
  binding constraint beyond this scale is identifiability rather than computation, since the
  number of GRN parameters grows as $n^2$. We now make explicit in the new paragraph that this
  ceiling takes the practical form of a trade-off between gene-set size and the amount of prior
  structure imposed on the edges: in the fully unconstrained setting we recommend not exceeding
  the order of 100 genes, whereas accepting more prior structure --- through a prior network
  $\theta^0$ or zeroed entries --- allows a correspondingly larger gene
  set for a comparable level of reliability.
\end{itemize}}

\revcomment{Accuracy and availability of kinetic parameters in realistic applications}{%
What level of uncertainty is typical for protein degradation-rate estimates in the systems
considered? Do literature-derived estimates generally fall within the uncertainty range
where the method remains robust?}
\response{%
We have modified the relevant paragraph in Section~2.3 connecting Figure~4F to realistic biological
settings. Protein half-lives in mammalian cells typically range 40--100h with a larger range up to 10--200~h for more extreme cases (Dear et al. 2024, \textit{PNAS}; Schwanh\"{a}usser et al. 2011,
\textit{Nature}). As an example of uncertainty, Rolfs et al. (2021, \textit{Nature Communications}) report $\pm$0.9d\ of median uncertainty in half-life in their mouse tissue atlas of 3106 proteins. Figure~4F shows CardamomOT maintains high AUPR
up to $\pm$50\% perturbation. Typical uncertainties correspond to roughly $\pm$50 \% relative error on the half-life for the most common range, i.e. at the boundary of the regime where accuracy is preserved; we therefore also report robustness up to $\pm$400 \%, and note that performance remains well above baseline throughout.}

\revcomment{Additional perturbation predictions for less-characterized genes}{%
To demonstrate discovery potential, could the authors provide a proof-of-concept showing
that CardamomOT can identify less-characterized genes with potentially important regulatory
roles?}
\response{%
We have added a supplementary table (Table~S3) listing the top candidate hub regulators by inferred outgoing regulatory strength for each experimental dataset, and updated the text in Section~2.3 and Section~2.6 accordingly.

We agree that the perturbations tested in silico in the original submission — \textit{Dnmt3a} and \textit{CHGA}, identified as high-confidence candidate regulators before any perturbation experiment, and \textit{Obox6} and \textit{Zfp42}, whose predicted effect on reprogramming efficiency matches published experimental validation~[15] — do not by themselves constitute a proof-of-concept for \emph{less-characterized} genes, since they were selected precisely because independent validation was already available. To address this directly, Table~S3 also surfaces regulators with little or no established role in the corresponding differentiation process, which we now discuss explicitly in Section~2.6: for example \textit{HAND2-AS1}, a long non-coding RNA antisense to \textit{HAND2}, in the Kameneva sympathoadrenal network, and \textit{Dmrtc2} and the poorly annotated \textit{Gm26917} in the Schiebinger reprogramming network. These candidates were not selected on the basis of prior functional knowledge, and their co-occurrence with experimentally validated regulators in the same ranked output illustrates the potential of CardamomOT to generate testable hypotheses about previously uncharacterized genes, in addition to recovering known regulatory relationships. As in the original submission, we frame all candidates listed in Table~S3 as testable hypotheses, and note explicitly that experimental validation of the newly proposed candidates is beyond the scope of this work.

We also note, as an illustration of the care required when interpreting such ranked outputs, that a subset of top-ranked candidates in the Kameneva dataset (e.g. \textit{RPL30}, \textit{ATP5F1E}) correspond to highly and uniformly expressed housekeeping genes rather than plausible regulators; we flag this explicitly in Section~2.3 so that ranked outputs of this kind are read alongside biological plausibility rather than taken at face value.}

\revcomment{Clarification of simulated-data noise model}{%
It is not fully clear whether simulations include technical noise resembling scRNA-seq
measurement noise. Could the authors clarify and, if not, add a robustness analysis?}
\response{%
The simulated data generated by HARISSA include \textit{intrinsic} biological stochasticity
(from the two-state model) but not \textit{technical} noise. We have added a supplementary robustness analysis (Figure 4G) introducing 7 increasing levels of binomial dropout imitating scRNA-seq noise. CardamomOT maintains high AUPR scores with slightly increasing variability at high levels of dropout, confirming robustness to moderate technical noise. 
We updated text accordingly in Section~2.3 and Section~2.1.
\seealso{R1.1.}}

\revcomment{Usability and documentation of the software package}{%
The documentation is currently difficult to navigate. The authors should consider hosting
structured online documentation (e.g.\ ReadTheDocs) with clear separation of installation,
tutorials, API documentation, and worked examples.}
\response{%
We have restructured the CardamomOT GitHub repository and created a ReadTheDocs
documentation site at \url{https://cardamomot.readthedocs.io} (link added to the Data and
Code Availability section). The documentation now includes:
\begin{itemize}
  \item Installation guide
  \item Quick-start tutorial reproducing the main results of the paper
  \item API reference (auto-generated from docstrings)
  \item Three worked examples corresponding to the three experimental datasets
\end{itemize}}

\subsect{Minor Comments}

\revcomment{Degradation missing from line~168 description}{%
Around line~168, mRNA and protein degradation should also be mentioned alongside
activation, synthesis, and production.}
\response{%
Revised to: ``[\ldots] the stochastic interplay between gene activation, mRNA synthesis,
protein production, and their respective degradation processes.''}

\revcomment{Bursty regime assumption: does it hold for all genes?}{%
Around line~180, is the bursty regime ($k_{\mathrm{on}} \ll k_{\mathrm{off}}$) expected to
hold for all genes? What fraction can reasonably violate it?}
\response{%
The bursty regime is well supported for the majority of mammalian genes (Suter et al.,
2011; Rodriguez \& Larson, 2020, already cited). Constitutively expressed genes may not
satisfy $k_{\mathrm{on}} \ll k_{\mathrm{off}}$; for these, the NB approximation remains
valid but the two-mode approximation is less appropriate. We have amended Section~1.1 to note this limitation and point to the softmax extension (multi-mode model) as a potential solution.}

\revcomment{NeuralODE recalibration step (line~302) is difficult to follow}{%
It would be easier to follow if more details on the Neural ODE part could be provided.}
\response{%
We thank the reviewer for pointing out this gap. We have substantially expanded
Section~1.2.3 to clarify that this step, previously described as NeuralODE recalibration and renamed "kinetic rate recalibration via differentiable ODE" for the sake of clarity, is not a generic black-box network but a
direct parametrization of the mechanistic protein field of Eq.~(5), with the converged
GRN $\theta^*$ frozen and only the protein degradation rates $d_1$ (plus a small per-gene
scale correction) treated as trainable. We now explicitly describe: (i) the model form,
integrated with a differentiable adaptive-step ODE solver (\texttt{torchdiffeq}); (ii) the
teacher-forcing training procedure, in which each consecutive timepoint pair is integrated
starting from the \emph{observed} reconstructed protein state rather than from the model's
own earlier prediction, preventing compounding integration error, following standard
practice for single-cell dynamical models (PRESCIENT~[45], STORIES~[46]); (iii) the explicit
loss function, combining reconstruction MSE with shrinkage of $d_1$ toward its literature
prior; and (iv) the companion recalibration of the mRNA degradation rate $d_0$, which we now
show is a closed-form, regularized method-of-moments estimator of $\tau = d_1/d_0$
derived from Campbell's theorem for the underlying bursty PDMP, rather than an unspecified
variance-fitting step. A pseudocode block (Algorithm~2) summarizing the procedure
has been added.
}


\revcomment{Reference needed for hub-regulator claim (line~460)}{%
The statement that sparse networks with hub regulators are a hallmark of biologically
realistic network structure should be supported by a reference.}
\response{%
We have added references supporting this observation (Ouma et al. 2018, \textit{Plos Computational Biology} and Sanguinetti et al. 2019, book chapter).}

\revcomment{Long section on experimental GRN inference without a main-text figure}{%
The section beginning around line~453 is long and contains several important results on
experimental GRN inference without a main-text figure. Consider shortening or moving a
supplementary figure into the main text.}
\response{%
We agree, and have restructured this material rather than moving one of the existing
supplementary figures into the main text. The reason is that the length of this section
came mainly from descriptive passages --- the gene selection criteria, the definition of
the virtual stimulus nodes, and a detailed walkthrough of their inferred targets --- rather
than from results that a figure would help the reader interpret. Promoting
Figure~S1--S3 to the main text would have added a large multi-panel figure per dataset
without shortening the text, and the quantitative results now most relevant to this section
(reproducibility across runs, ranked hub regulators) are already reported in
Figure~S9 and Table~S3, both added in response to Reviewer~1.

We have therefore moved the descriptive material to a new
Appendix C (``Details on GRN inference for experimental datasets''),
which collects the gene selection procedure, the definition and role of the stimulus nodes
in each dataset, and the consistency check between their inferred targets and the known
biology of each protocol. The main text now retains only the results themselves: network
sparsity and hub structure, reproducibility across independent runs, the identity of the
top regulators and their agreement with the original studies, and the non-negligible
diagonal coefficients. This substantially reduces the section while keeping every
result, and each condensed statement points to the appendix for the corresponding detail.

We also note that the inferred experimental GRNs are already represented in the main text, in the left panels of Figure 7, which show the regulatory neighbourhood of each perturbed gene in the corresponding dataset; Section 2.3 is thus not without a main-text visual counterpart, and the section it feeds into is figure-supported.
}


\revcomment{Figure~6 panels J--O: legend missing}{%
Panels J--O would be easier to interpret if a legend were included directly in the figure.}
\response{%
A legend has been added directly within panels J--O of Figure~6, identifying reference
(gray), NB mixture (blue), and simulation (red) traces.}

\revcomment{Several figures: text too small}{%
Several figures would benefit from larger text, especially plot titles. For example, the
titles in Figure~6J--O are difficult to read at the current size.}
\response{%
We have increased font sizes for plot titles and axis labels throughout the figures to
ensure readability at print size. Figures have been re-exported.}

\revcomment{Non-convex objective: guidance on multiple initializations}{%
Around line~761, the authors note that the objective is non-convex. Should users run
CardamomOT with multiple initializations? Please provide guidance.}
\response{%
We have added practical guidance in Section~3 and in the online documentation:
\begin{enumerate}
  \item For networks up to $\sim$20 genes with good temporal resolution, a single run is
        typically sufficient (Figure~4A--E shows convergence to similar GRN structures
        regardless of initialization).
  \item For larger or sparser networks, we recommend 3--5 runs with different random seeds;
        high variability in edge weights across runs signals poor identifiability for that
        edge.
\end{enumerate}
\seealso{R1.2.}}

% -------------------------------------------------------------

\noindent We hope that the revised manuscript and this response letter satisfactorily
address all concerns raised by the Editor and Reviewers. We remain available to provide
any additional clarifications or analyses.

\bigskip
\noindent Sincerely,

\noindent Yann Maugé and Elias Ventre

\end{document}